% Author continuation: S tends to infinity only at fixed m,R.
\section{The exact jump spectrum of the full marked physical roof}
\label{sec:js-spectrum}

The absence of a singular-continuous derivative in
Section~\ref{sec:sbv-roof} identifies the entire nondecaying
mean-square part of the differentiated physical roof transform.
The result keeps the two displacement frequencies and the
occupation frequency; it does not average over collision counts
or change a conditioning event. It is an exact identity at
every finite count, with a collision-uniform bound and a necessary
consequence of the original raw limit. Its high-frequency limit
is not interchanged with the collision limit.

Use the source and tests of \eqref{eq:wo-measures}. Write
$\ell=(k_1,k_2,n)$, $f_\ell=f^{w;a,b,\Psi}_{(n,k),m,R}$,
$D=\mathcal V_m^w(a,b,\Psi)$ and $A_V=a_0b_0$. At fixed $m,R$
only finitely many labels are attainable. By
Theorem~\ref{thm:sbv-physical},
\begin{equation}\label{eq:js-decomposition}
 Df_\ell=d_\ell(t)\,dt+
       \sum_{u\in\mathcal S_{m,R}}a_\ell(u)\delta_u,\qquad
 a_\ell(u)=f_\ell(u+)-f_\ell(u-),
\end{equation}
where $d_\ell\in L^1$ and $\mathcal S_{m,R}$ is finite.
Zero coefficients are permitted. The geometric support is not
asserted uniformly separated. All coincident contributions to
a given $(\ell,u)$ have already been combined with their signs.
In particular
\begin{equation}\label{eq:js-budget}
 \sum_\ell\left(\|d_\ell\|_1+\sum_u|a_\ell(u)|\right)\le D.
\end{equation}
This is the existing complete all-label variation budget applied
to the mutually singular parts of its exact derivative.

Use the full marked transform
$\mathcal F(\vartheta,s)=\mathcal F_{m,R}^{w,V}(\vartheta,s)$
of \eqref{eq:ui-transform}, with
$\vartheta\in[-\pi,\pi]^3$. For $S>0$ define
\begin{align}
 \mathcal A(\vartheta,s)
 &=\sum_{\ell,u}e^{i\vartheta\cdot\ell+isu}a_\ell(u),&
 \mathcal G(\vartheta,s)
 &=\sum_\ell e^{i\vartheta\cdot\ell}\widehat d_\ell(s),
                                                   \label{eq:js-transforms}\\
 \mathcal J_{m,R}(w,V)&=\sum_{\ell,u}|a_\ell(u)|^2,&
 \mathcal W_S&=\frac1{2S(2\pi)^3}
 \int_{-S}^S\int_{[-\pi,\pi]^3}
       s^2|\mathcal F(\vartheta,s)|^2\,d\vartheta\,ds.
                                                   \label{eq:js-energy}
\end{align}
The convention is $\widehat f(s)=\int e^{ist}f(t)\,dt$.

\begin{theorem}[All-label physical jump spectrum]
\label{thm:js-spectrum}
For every actual $m\ge1$, $R\in[R_-,R_+]$, original specified
weight and tests,
\begin{equation}\label{eq:js-identity}
 -is\mathcal F(\vartheta,s)
       =\mathcal A(\vartheta,s)+\mathcal G(\vartheta,s),
 \qquad
 \lim_{S\to\infty}\mathcal W_S
       =\mathcal J_{m,R}(w,V)\le D^2.
\end{equation}
The first identity holds for every real $s$ and full-torus
frequency. At fixed $m,R$,
\begin{equation}\label{eq:js-asymptotic}
 \sup_\vartheta
 |s\mathcal F(\vartheta,s)-i\mathcal A(\vartheta,s)|
       \longrightarrow0\qquad(|s|\to\infty).
\end{equation}
For a precise finite-frequency error, put
\begin{align}
 \mathcal Q_S&=\frac1{2S}\int_{-S}^S
                         \sum_\ell|\widehat d_\ell(s)|^2\,ds,
                                                   \label{eq:js-q}\\
 \mathcal C_S&=\sum_\ell
 \sum_{\substack{u,v\in\mathcal S_{m,R}\\u\ne v}}
 |a_\ell(u)a_\ell(v)|
       \min\left\{1,\frac1{S|u-v|}\right\}.
                                                   \label{eq:js-cross}
\end{align}
Then
\begin{equation}\label{eq:js-finite}
 |\mathcal W_S-\mathcal J_{m,R}(w,V)|
       \le\mathcal C_S+2D\sqrt{\mathcal Q_S}+\mathcal Q_S.
\end{equation}
Both $\mathcal Q_S$ and $\mathcal C_S$ tend to zero for each
fixed $m,R$. No positive uniform lower bound for distinct
roof gaps, or rate uniform in $m,R$, is asserted.
\end{theorem}
\begin{proof}
Compact support and distributional integration by parts give
\[
 -is\widehat f_\ell(s)=\widehat d_\ell(s)
                      +\sum_u a_\ell(u)e^{isu}.
\]
At $s=0$ this is the vanishing total mass of a compactly
supported derivative. Summing the finite label set gives the
first identity, with the displayed sign.

For $g\in L^1(\mathbb R)$, approximate $g$ in $L^1$ by finite
linear combinations of bounded interval indicators. Their
transforms tend to zero by direct integration, and the transform
error is at most the $L^1$ error, uniformly in frequency.
Thus $\widehat g(s)\to0$. Apply this to every $d_\ell$.
The finite label sum yields
$\sup_\vartheta|\mathcal G(\vartheta,s)|
 \le\sum_\ell|\widehat d_\ell(s)|\to0$,
which proves \eqref{eq:js-asymptotic}.

Full-torus orthogonality, before the roof-frequency average,
gives exactly
\begin{equation}\label{eq:js-parseval}
 \mathcal W_S=\frac1{2S}\int_{-S}^S\sum_\ell
 \left|\widehat d_\ell(s)+\sum_u a_\ell(u)e^{isu}\right|^2ds.
\end{equation}
No factor counting labels is introduced. The purely atomic
part has mean
\begin{equation}\label{eq:js-sinc}
 \mathcal J_{m,R}(w,V)+
 \sum_\ell\sum_{u\ne v}a_\ell(u)\overline{a_\ell(v)}
                 \frac{\sin(S(u-v))}{S(u-v)}.
\end{equation}
The second term has modulus at most $\mathcal C_S$.
The diagonal is the square of the combined coefficient at each
coincident roof, not the sum of squares of an artificial partition.

At every $s$,
\[
 \sum_\ell\left|\sum_u a_\ell(u)e^{isu}\right|^2
 \le\left(\sum_{\ell,u}|a_\ell(u)|\right)^2\le D^2.
\]
Cauchy--Schwarz in the frequency average bounds the mixed
term in \eqref{eq:js-parseval} by $2D\sqrt{\mathcal Q_S}$;
the remaining term is $\mathcal Q_S$. This proves
\eqref{eq:js-finite}. The finite cross sum tends to zero.
The function $\sum_\ell|\widehat d_\ell(s)|^2$ is bounded by
$D^2$ and tends to zero at both ends of the line, so its
symmetric averages tend to zero. The claimed limit and upper
bound follow. All operations occur at the specified finite
count and radius, after including the entire physical source.
\end{proof}

\begin{remark}[Coalescing roofs and the limit order]
\label{rem:js-coalescence}
The cross term \eqref{eq:js-sinc} retains coalescing roof values.
For example, the derivative $\delta_0-\delta_\delta$ of
$\mathbf1_{(0,\delta)}$ has atomic energy $2$ for every
$\delta>0$, but finite-frequency energy
$2-2\sin(S\delta)/(S\delta)$, tending to zero as
$\delta\downarrow0$ with $S$ fixed. This analytical check is
not a substitute billiard model. A uniform variation bound
alone does not justify moving the frequency limit through a
radius supremum or collision limit. Nor does the
$o(1/|s|)$ remainder assert absolute integrability.
\end{remark}

\subsection{A necessary central spectrum at the original normalization}

Return to $p_\ell=p_{n,R}(k,m,\cdot)$ and
$D_0=D_{\rm all}^{1}$. Fix $M_1>M>0$ and a $C^1$ function
$0\le\psi\le1$ on $[0,\infty)$ equal to one on $[0,M^2]$
and zero on $[M_1^2,\infty)$. Define the analytical roof test
\begin{equation}\label{eq:js-central-cutoff}
 \chi_{\ell,m,R}(t)=\psi\left(
 \left|\frac{(k_1,k_2,t,n)-m\mu_R}{\sqrt m}\right|^2\right),
 \qquad h_\ell(t)=\chi_{\ell,m,R}(t)p_\ell(t).
\end{equation}
This is a bounded observation of the original physical record;
it does not replace its entrance law, terminal test or
microscopic interval. Its choice is fixed before the collision
limit. With the exact jumps $a_\ell(u)$ of $p_\ell$, put
\[
 \mathcal J^\chi_{m,R}
     =\sum_{\ell,u}|\chi_{\ell,m,R}(u)a_\ell(u)|^2.
\]
The same Fourier calculation as in the proof of
Theorem~\ref{thm:js-spectrum} gives the energy limit
$\mathcal J^\chi_{m,R}$ for the marked transform of $(h_\ell)$.
This is justified by its following product decomposition, rather
than by classifying the label-dependent cutoff as an endpoint
test in \eqref{eq:wo-measures}. The BV product rule gives
\[
 D h_\ell=(\chi_{\ell,m,R}d_\ell+
                 \chi'_{\ell,m,R}p_\ell)\,dt
       +\sum_u\chi_{\ell,m,R}(u)a_\ell(u)\delta_u.
\]
Its summed variation is at most
$D_0+2M_1\|\psi'\|_\infty/\sqrt m$, since
$|\chi'|\le2M_1\|\psi'\|_\infty/\sqrt m$ and
$\sum_\ell\|p_\ell\|_1\le1$.

\begin{theorem}[The raw target forces a complete normalized jump spectrum]
\label{thm:js-raw-necessary}
Let $E_m(M_1)$ be the raw error in \eqref{eq:v76-target}
at central parameter $M_1$, before taking its limit.
For every $m$,
\begin{equation}\label{eq:js-necessary}
 \sup_R m^2\mathcal J^\chi_{m,R}\le2D_0E_m(M_1).
\end{equation}
Consequently the original raw target implies
\begin{equation}\label{eq:js-necessary-limit}
 \lim_{m\to\infty}\sup_R m^2\mathcal J^\chi_{m,R}=0.
\end{equation}
This is a necessary condition, not a sufficient one.
\end{theorem}
\begin{proof}
If $\chi_{\ell,m,R}(u)\ne0$, $u$ is interior to the
larger central roof interval. The retained finite-peak reference
is continuous there, including all near-unit branches.
Approach $u$ from both sides by Lebesgue points in that interval.
The definition of $E_m(M_1)$ and the BV traces give
\[
 |m^2p_\ell(u\pm)-\mathcal L_{m,R}(k_1,k_2,u,n)|
                       \le E_m(M_1).
\]
Hence $m^2|a_\ell(u)|\le2E_m(M_1)$, equivalently the
physical trace argument of Proposition~\ref{prop:sbv-central-jumps}.
Multiply by $|a_\ell(u)|\chi_{\ell,m,R}(u)^2$ and sum
all labels and roofs. The complete budget
$\sum_{\ell,u}|a_\ell(u)|\le D_0$ proves
\eqref{eq:js-necessary}. The target at $M_1$ proves the limit.
The allowable null set may depend on $R$. The frequency limit
identifying $\mathcal J^\chi$ was taken at fixed $m,R$;
it has not been passed through a supremum or collision limit.
\end{proof}

\subsection{The inherited zero-clearance payment in this signed spectrum}

Take the one-copy union of positive-visited critical germs with
at least one zero-clearance unvisited-disk witness covered by
Theorem~\ref{pc:fan}. Unite actual incident sectors first.
When a germ has several witnesses, assign it once by a fixed
ordering of its finite flight indices, nearby centers and signs;
then split by its exact-label sectors. These sectors partition
each physical fan, so the sum over labels introduces no extra
multiplicity. This is an enumeration
of the actual source, not a change of the source or event.
At each germ include its actual sector and continuous bounded
test value. Let $b_\ell(u)$ be this part of the atomic coefficient
in \eqref{eq:js-decomposition}, grouped at the same exact label
and roof. For the positive unweighted source these incoming
coefficients are nonnegative; tested coefficients can be complex.
They need not be the entire atom at that roof.

\begin{corollary}[Paid fans with every roof coincidence retained]
\label{cor:js-paid-fans}
With $\sigma=47/59$ and the inherited $C_{\rm pc}$, put
$B_m=A_VC_{\rm pc}m\sigma^{(m-1)/2}$. Then
\begin{equation}\label{eq:js-paid-fans}
 \begin{split}
 \sum_{\ell,u}|b_\ell(u)|&\le B_m,\\
 \left|\sum_{\ell,u}|a_\ell(u)|^2
       -\sum_{\ell,u}|a_\ell(u)-b_\ell(u)|^2\right|
                         &\le2DB_m+B_m^2.
 \end{split}
\end{equation}
For fixed original guard width and fixed test norms, the last
bound vanishes after multiplication by any prescribed power of
$m$, uniformly in $R$. The same conclusion holds with the
central tests \eqref{eq:js-central-cutoff}.
This transfers the already proved zero-clearance fan payment
to the new signed spectrum. It makes no claim about the
positive-clearance critical sources, other constant-roof side
charges, or the absolutely continuous derivative.
\end{corollary}
\begin{proof}
The positive fan bound \eqref{pc:fan-bound} includes both
impact signs, every actual marked disk and visited word, and
the original finite section, record and roof restrictions.
The once-only union is bounded by the positive sum over all
witnesses. Sum over $j=0,\ldots,m-1$ and use
$\max(j,m-1-j)\ge(m-1)/2$. This gives
$C_{\rm pc}m\sigma^{(m-1)/2}$. The guard is at most one,
and the continuous test multiplier has modulus at most $A_V$
at the fixed positive-visited germ. Combining true labels and
coincident roofs before applying triangle inequalities gives
the first bound.

For arbitrary complex vectors,
\[
 \big|\|a\|_2^2-\|a-b\|_2^2\big|
       \le2\|a\|_2\|b\|_2+\|b\|_2^2.
\]
Use $\|a\|_2\le\|a\|_1\le D$ and
$\|b\|_2\le\|b\|_1\le B_m$. Multiplication of coefficients
by $0\le\chi\le1$ decreases these norm bounds. Exponential
decay dominates every fixed power of $m$. The conditioning
is already in $D$ and $C_{\rm pc}$; there is no second
division by $c_*$. The proof invokes the established
physical fan theorem and does not claim that theorem as new.
\end{proof}

\paragraph{Relation to the unchanged raw objective.}
The special-BV theorem excludes a Cantor derivative for the
specified physical observations. The exact all-label spectrum
identifies its atomic part with the full signed cancellations
and arithmetic labels. These results do not establish vanishing
of the remaining normalized jumps or continuous local derivative
charges. Smooth narrow spikes can have no jumps at all.
The original same-roof shortage \eqref{eq:lc-final-shortage}
remains the closing condition on its existing route.
The alternative sufficient bridge of
Proposition~\ref{prop:ui-inherited-bridge} still requires its
complete normalized local derivative charge, in the original
order of limits. Fixed-count structure and the new frequency
identity are not uniform dynamical estimates of that charge.
Microscopic conditioning and the four-coordinate inverse
retain their actual denominator and Jacobian hypotheses.

